\section{Limitations}
\label{sec:limitations}

The analysis assumes that candidate answers can be mapped to a finite set and
that generations are conditionally independent given a question and model.
Correlated decoding traces or unstable answer extraction would require a
different concentration radius.  The uniform radius is also conservative: it
pays a logarithmic factor for all observed model--answer--question
coordinates, even when many answers are rare or impossible.  Data-dependent
or variance-sensitive buffers may be substantially smaller.

The main guarantee targets the clean population risk of weights fitted from
finite-generation training margins.  A deployment that estimates answer
frequencies from fresh finite rollouts has a second perturbation.  The retained
fixed-weight sign bound describes that effect, but controlling it uniformly at
the random fitted weight would need an additional margin condition or a
selective prediction rule.

The optional fast \(Q\)-rate uses identifiable risk growth and global
separation.  These assumptions are intentionally stronger than a margin
condition and should be checked rather than treated as generic properties of
LLM benchmarks.  Finally, the tracked real-data analysis predates the
calibrated buffered estimator.  It supports the finite-sample motivation but
does not replace the direct \(Q\)-by-\(N\) evaluation specified in
Section~\ref{sec:experiments-required}.

\section{Conclusion}
\label{sec:conclusion}

Finite-generation ensemble weighting is governed by two sampling operations
that should not be collapsed into one.  Questions determine how well the
weight vector generalizes; generations determine whether a question's
gold-versus-contender margin can be resolved.  Buffering at the latter
uncertainty scale yields a clean comparison with the best fixed population
weight and a robust \(Q^{-1/2}+N^{-\beta/2}\) guarantee up to logarithmic
factors.  The hard-margin counterexample also draws a useful boundary: local
separation from zero is enough to control finite-generation errors, but not
enough to identify the best region of the weight simplex quickly.  Faster
learning across questions requires additional risk geometry.

This theorem hierarchy leaves a concrete empirical next step.  The existing
MILP already exposes the required margin threshold; a direct benchmark study
must now vary \(Q\) and \(N\), compare buffered and unbuffered fits, and retain
the split-level artifacts needed to distinguish the two terms.  Until then,
the theorem, the obstruction, and the legacy diagnostics make separate,
auditable claims.
